\documentclass{article}
\usepackage[margin=2cm]{geometry}
\usepackage{amsmath}
\usepackage{graphicx}
\usepackage{booktabs}
\usepackage[utf8]{inputenc}
\usepackage{ragged2e}
\setlength{\emergencystretch}{3em}
\usepackage{paracol}
\begin{document}
\sloppy

\begin{paracol}{2}
\RaggedRight

components in I. Let 0*(o) e O the initial parameter.
Let 0,0w=0o+VJ0*wo) and0*)=
0o + VJ\. (0* (o)). Then, under Assumption, we have:


\begin{equation}
\begin{array} { r l } & { * ~ }\end{array}
\end{equation}


Thus, we maximize the L2norm of the gradient com-
Since we have at our disposal only samples D collected with
the current policy e*() and in the current environment wo,
we have to perform an offdistribution optimization over w.
To this end, we employ an approach analogous to that of
where we optimize the empirical version of the objective
with a penalization that accounts for the distance between
the distribution over trajectories:


\begin{equation}
\begin{array} { r } { * } \\ { C _ { I } ( \omega / \omega _ { 0 } ) = \left \| \widehat { \nabla } _ { \theta } J _ { M _ { \mathrm { c o t e n } } } ( \theta ^ { * } ( \omega _ { 0 } ) ) | _ { \overline { I } } \right \| _ { 2 } ^ { 2 } - \zeta \sqrt { \frac { \hat { d } _ { 2 } ( \omega \| \omega _ { 0 } ) } { n } }, ~ ~ ( 1 ) } \\ { \underset { \mathrm { a l h e s i m } } { \widehat { \psi } } \mathrm { ~ \textit ~ { ~ r ~ } ~}}\end{array}
\end{equation}


where VJM/ (0* (o)) is an off-distribution estimator of
the gradient VJM\. (6*(o) using samples collected with
wo, d2 is the estimated 2-Renyi divergence that works as
a penalization to discourage too large updates and (  0
is a regularization parameter. The expression of the estimated gradient, 2-Renyi divergence and the pseudocode are
reported in Appendix.


\subsection*{Experimental Evaluation}


In this section, we present the experimental evaluation of the
identification rules in three RL domains. To set the values of
c(l) we resort to the Wilks asymptotic approximation (Theorem) to enforce (asymptotic) guarantees on the type I error. For Identification Rule we perform 2d statistical tests by
using the same dataset D. Thus, we partition \& using Bonis the -quintile of a chi square distribution with I degrees
of freedom. Instead, for Identification Rule, we perform d


\subsection*{Discrete Grid World}


The grid world environment is a simple representation of a
two-dimensional world (5  5 cells) in which an agent has to
reach a target position by moving in the four directions. The
goal of this set of experiments is to show the advantages
icy space identification using rule. The initial position of the
agent and the target position are drawn at the beginning of
each episode from a Boltzmann distribution . The agent
plays a Boltzmann linear policy  with binary features  indicating its current row and column and the row and column
of the goal. For each run, the agent can control a subset I*
of the parameters 0 associated with those features, which
is randomly selected. Furthermore, the supervisor can configure the environment by changing the parameters w of the
initial state distribution . Thus, the supervisor can induce
the agent to explore certain regions of the grid world and,
consequently, change the relevance of the corresponding parameters in the optimal policy.

\switchcolumn
\RaggedRight

Figure shows the empirical  and , i.e. the fraction of
parameters that the agent does not control that are wrongly
selected and the fraction of those the agent controls that are
not selected respectively, as a function of the number n of
episodes used to perform the identification. We compare two
cases: conf where the identification is carried out by also
configuring the environment, i.e. optimizing Equation (1),
and no-conf in which the identification is performed in the
original environment only. In both cases, we can see that
 is almost independent of the number of samples, as it is
directly controlled by the critical value c(1). Differently, 
decreases as the number of samples increases, i.e. the power
of the test 1   increases with n. Remarkably, we observe
that configuring the environment gives a significant advantage in understanding the parameters controlled by the agent
w.r.t using a fixed environment, as  decreases faster in the
conf case. This phenomenon also justifies empirically our
choice of objective (Equation (1)) for selecting the new environment. Hyperparameters, further experimental results, together with experiments on a continuous version of the grid
world, are reported in Appendix.


\subsection*{Minigolf}


In the Minigolf environment, an agent hits a ball using a putter with the goal of reaching the hole in the minimum number of attempts. Surpassing the hole causes the termination
of the episode and a large penalization. The agent selects the
force applied to the putter by playing a Gaussian policy linear in some polynomial features (complying to Lemma) of
the distance from the hole () and the friction of the green
(f). We consider two agents: / has access to both the  and
f whereas s/2 knows only x. Thus, we expect that / learns
a policy that allows reaching the hole in a smaller number of
hits, compared to /2, as it can calibrate force according to
friction; whereas s/2 has to be more conservative, being unfor the two agents, the best putter length w, i.e. the configurable parameter of the environment. In this experiment, we
want to highlight that knowing the policy space might be of
crucial importance when learning in a ConfMDP.


Figure-left shows the performance of the optimal policy
as a function of the putter length w. We can see that for agent
 the optimal putter length is w - 5 while for agent
= 11.5. Figure-right compares the performance
of the optimal policy of agent / when the putter length w
is chosen by the supervisor using four different strategies.
In (i) the configuration is sampled uniformly in the interval
[1, [5]. In (ii) the supervisor employs the optimal configuration for agent /1 ( = 5), i.e. assuming the agent is aware
of the friction. (iii) is obtained by selecting the optimal configuration of the policy space produced by using our identification rule. Finally, (iv) is derived by employing an oracle
that knows the true agent's policy space ( = 11.5). We can
see that the performance of the identification procedure (i)
is comparable with that of the oracle (iv) and notably higher
than the performance when employing an incorrect policy
space (ii). Hyperparameters and additional experiments are

\end{paracol}

\end{document}
